\section{Finite-cell realization and comparison of experiments}
\label{sec:v33-moments}

The realization algebra has a wider scope than the collision theorem.
It is useful to state it at that scope before returning to the monomial
geometry. The static coefficient set is a finite-output postprocessing
region in the sense of Blackwell \cite{Blackwell1953}. Its iteration
records the joint law of the retained label and the latent parameter.
Neither this postprocessing interpretation nor polynomial propagation
asserts the attainable rank or the collision profile proved above.

Let $\Theta$ be a standard Borel space, let $\mu$ be any probability
on $\Theta$, and let $k_1^a,\ldots,k_J^a:\Theta\to[0,1]$ be Borel
functions with $\sum_k k_k^a=1$. Here $a$ denotes a fixed known
specification of the cells. No monomial structure, density of $\mu$,
or strict positivity of the cells is required in this section until
we explicitly return to collision limits. Conditional on $t\in\Theta$,
the raw cells on successive trials are independent with law $k^a(t)$.
The command lies
in $G=[\eta,1-\eta]^J$, where $0<\eta<1/2$; raw cell $k$ is
reported as $k$ with probability $g_k$ and otherwise as $\dagger$.
The command at time $n$ has prescribed law $\nu_n$, independently of
the latent parameter, all previous variables and the fresh private coins.
The controller reads only its preceding label, current command and
current report, and retains one of $M$ labels after each report.

At checkpoint $n<N$ there is a finite query menu with fixed
nonnegative weights $\omega_{nq}$ summing to one. Each future
likelihood $H^a_{nq}\in[0,1]$ is the product of at most $N-n$
cell-linear likelihood factors. In particular it is a polynomial of
degree at most $N-n$ in the cells. The query is selected independently
after retention. Mean regret is the squared difference from the
full-history prediction, averaged over this prescribed experiment,
and the common-controller criterion is the maximum of these means.
The notation of Definition~\ref{def:monomial-risks-v24} is extended
to this finite-cell model with exactly those quantifiers. The positive
monomial model and its uniform command laws are a specialization.
A history of zero evidence contributes zero to every integral; its
posterior prediction can be assigned any value in $[0,1]$.

\subsection{The implementability body}
For a command law $\nu$, put $m_k=\int_G g_k\,\nu(\mathrm dg)$ and define
\begin{equation}\label{eq:v33-allocation}
 \mathscr Z_M(\nu)=\left\{z\in\R^{J\times M}:
 z_{kj}=\int_G(1-g_k)u_j(g)\,\nu(\mathrm dg),\quad
 u:G\longrightarrow\Delta_M\ \hbox{Borel}\right\},
\end{equation}
where $\Delta_M$ is the probability simplex. Define the one-state
implementability body by
\begin{equation}\label{eq:v33-body}
 \mathscr A_M(\nu)=
 \left\{ A\in\R^{J\times M}: A_{kj}=m_k r_{kj}+z_{kj},\quad
       r_k\in\Delta_M\ (1\le k\le J),\quad z\in\mathscr Z_M(\nu)\right\}.
\end{equation}
Every row of $A$ sums to one. The shared array $z$ is essential: all
raw cells that produce failure must use the same failure update.


\begin{proposition}[The finite-output Blackwell region]
\label{prop:v34-blackwell}
On $\Omega=G\times\{1,\ldots,J,\dagger\}$ define the experiment
with auxiliary index $k\in\{1,\ldots,J\}$ by
\begin{equation}\label{eq:v34-experiment}
 \mathsf E_k(\mathrm dg,\mathrm dx)=\nu(\mathrm dg)
       \{g_k\delta_k(\mathrm dx)+(1-g_k)\delta_\dagger(\mathrm dx)\}.
\end{equation}
Then, exactly,
\begin{equation}\label{eq:v34-blackwell}
 \mathscr A_M(\nu)=\{\mathsf E U:
       U\text{ is a Borel Markov kernel from }\Omega\text{ to }[M]\}.
\end{equation}
The actual observation law given $t$ is the mixture
$\sum_k k_k^a(t)\mathsf E_k$. The index of this auxiliary experiment
is not an additional observation supplied to the controller.

For a prior $\pi$ on the auxiliary indices and a reward table
$r(k,j)$, its optimal finite-action Bayes reward is the support function
of $\mathscr A_M(\nu)$ at $C_{kj}=\pi_k r(k,j)$. Moreover, for
$k\ne l$ the measure
\[
 \lambda_{kl}(\mathrm dg,\mathrm dx)
 =\min(1-g_k,1-g_l)\nu(\mathrm dg)\delta_\dagger(\mathrm dx)
\]
is a common submeasure of $\mathsf E_k$ and $\mathsf E_l$.
Its mass is preserved by every postprocessing, and hence
\begin{equation}\label{eq:v34-blackwell-overlap}
 \operatorname{TV}(A_k,A_l)
 \le 1-\lambda_{kl}(\Omega)\qquad(A\in\mathscr A_M(\nu)),
\end{equation}
where total variation between probability rows is half their
$\ell^1$ distance.
\end{proposition}
\begin{proof}
Direct integration gives
\[
 (\mathsf E U)_{kj}
   =\int_G\{g_k U(j\mid g,k)
                  +(1-g_k)U(j\mid g,\dagger)\}\,\mathrm d\nu.
\]
The accepted term is $m_k r_{kj}$ for the probability row obtained
by $g_k$-weighted averaging of the accepted kernel. The failure terms
all use the same simplex-valued function $U(\cdot\mid g,\dagger)$,
so belong to $\mathscr Z_M(\nu)$. Conversely, realize any
$m_k r_{kj}+z_{kj}$ by the constant accepted kernel $r_k$ and a
common failure kernel representing $z$. This proves
\eqref{eq:v34-blackwell}. Mixing the conditional raw-cell laws gives
the asserted physical observation law.

For any postprocessing the Bayes reward is
$\sum_{k,j}\pi_k r(k,j)(\mathsf E U)_{kj}$; taking the supremum
proves the support interpretation. Finally $\lambda_{kl}\le\mathsf E_k$
and $\lambda_{kl}\le\mathsf E_l$, so the row vector $\lambda_{kl}U$
is dominated by both output rows. Its entries sum to
$\lambda_{kl}(\Omega)$ because $U$ is a Markov kernel. Thus
$\sum_j\min(A_{kj},A_{lj})\ge\lambda_{kl}(\Omega)$, which is
\eqref{eq:v34-blackwell-overlap}.
\end{proof}

This identifies the classical framework of the body. The next lemma
computes its support for the acceptance/failure channel and proves the
compactness needed for its finite-horizon iteration. In particular,
the support calculation is a Bayes decision calculation, not a separate
principle of comparison of experiments.

\begin{lemma}[One-step implementability]\label{lem:v33-body}
The set $\mathscr A_M(\nu)$ is compact and convex. A matrix $A$ belongs
to it if and only if there are Borel transition probabilities
$U(j\mid g,x)$ such that
\begin{equation}\label{eq:v33-realize-row}
 A_{kj}=\int_G\bigl[g_k U(j\mid g,k)
                 +(1-g_k)U(j\mid g,\dagger)\bigr]\,\nu(\mathrm dg).
\end{equation}
For every $C\in\R^{J\times M}$ its support function is
\begin{equation}\label{eq:v33-support}
 \begin{split}
 \sup_{A\in\mathscr A_M(\nu)}\sum_{k,j}C_{kj}A_{kj}
 &=\sum_k m_k\max_j C_{kj}\\
 &\quad+\int_G\max_j\sum_k(1-g_k)C_{kj}\,\nu(\mathrm dg).
 \end{split}
\end{equation}
The body and this formula do not involve the calibration or the prior.
\end{lemma}
\begin{proof}
The measurable simplex-valued functions form a weak-star compact convex
subset of $L^\infty(\nu)^M$. Indeed, nonnegativity and the equation that
the coordinates sum to one are weak-star closed, and the coordinates
are bounded by one. Integration against the finitely many functions
$1-g_k$ is weak-star continuous. Consequently $\mathscr Z_M(\nu)$,
and hence $\mathscr A_M(\nu)$, are compact and convex. Equivalence
classes have Borel representatives on the command cube; null-set
values can be chosen in the simplex.

For a transition kernel in \eqref{eq:v33-realize-row}, its accepted
part has the form $m_k r_{kj}$, where
$r_{kj}=m_k^{-1}\int g_k U(j\mid g,k)\,\mathrm d\nu$. Here $m_k>0$.
Its failure part belongs to \eqref{eq:v33-allocation}. Conversely,
realize $m_k r_{kj}$ by the constant accepted-report rule $r_k$ and
realize $z$ by a kernel occurring in its definition. This proves both
directions without imposing separate failure rules for different raw
cells. Finally the accepted terms in a linear functional are maximized
at simplex vertices. For the failure term choose the smallest index
maximizing the finitely many functions
$\sum_k(1-g_k)C_{kj}$. This is a Borel rule, attains the pointwise
maximum, and proves \eqref{eq:v33-support}.
\end{proof}

\begin{corollary}[A constraint imposed by the common failure report]
\label{cor:v33-overlap}
For every $A\in\mathscr A_M(\nu)$ and any raw cells $k,l$,
\[
 \sum_j\min(A_{kj},A_{lj})
 \ge\int_G\min(1-g_k,1-g_l)\,\nu(\mathrm dg)\ge\eta.
\]
When $J,M\ge2$, the implementability body is a proper subset of the
product of stochastic row simplices.
\end{corollary}
\begin{proof}
Use the common failure kernel $u$ in \eqref{eq:v33-allocation}.
Both $A_{kj}$ and $A_{lj}$ dominate
$\int\min(1-g_k,1-g_l)u_j(g)\,\mathrm d\nu$.
Sum over $j$ and use $\sum_j u_j=1$. Two disjoint point-mass rows
violate the resulting positive overlap, proving proper inclusion.
\end{proof}

\subsection{Occupancy polynomials and exact mean risk}
Write $T=N-1$. The initial label is fixed: independent private
randomization of it can be averaged into the first transition, since
no prediction is made before the first report. For each time $n\le T$ and each preceding state $i$, take
one matrix $A_n(i)\in\mathscr A_M(\nu_n)$. Starting from
$q_{0,i}=\mathbf1_{\{i=1\}}$, define
\begin{equation}\label{eq:v33-q-recursion}
 q_{n,j}(t)=\sum_{i=1}^M\sum_{k=1}^J
        q_{n-1,i}(t)k_k^a(t)A_{n,kj}(i).
\end{equation}
Let $H^a_{nq}$ be the future likelihood of query $q$, with weight
$\omega_{nq}$, and set
\begin{equation}\label{eq:v33-coefficients}
 \begin{split}
 w_{ni}&=\mu(q_{n,i}),\qquad
 b_{niq}=\mu(q_{n,i}H^a_{nq}),\\
 B_n(a)&=\int_{\mathcal H_n}\sum_q\omega_{nq}
                    p^a_{nq}(h)^2\,\mathbb P_n^a(\mathrm dh).
 \end{split}
\end{equation}
The offset $B_n$ is determined by the prescribed experiment, not by
the controller. In the principal uniform menu
$\omega_{nq}=1/s_n^{\rm q}$.

\begin{theorem}[Finite-cell realization of a common controller]
\label{thm:v33-moment-program}
For the finite-cell model just specified, with arbitrary $\mu$ and
prescribed independent command laws,
\begin{equation}\label{eq:v33-exact-value}
 \mathfrak R_{M,N}^{a,\mathrm{av}}
  =\min_{\substack{A_n(i)\in\mathscr A_M(\nu_n)\\1\le n<N,\ 1\le i\le M}}
       \max_{1\le n<N}
       \left\{B_n(a)-\sum_{i,q}\omega_{nq}
                        \frac{b_{niq}^2}{w_{ni}}\right\}.
\end{equation}
A quotient with $w_{ni}=0$ is defined to be zero. The minimum is
attained and is realized by one common $M$-label controller. Its decoder
is $d_{nq}(i)=b_{niq}/w_{ni}$ whenever $w_{ni}>0$.

More precisely, the polynomials in \eqref{eq:v33-q-recursion} are
exactly the functions
$\mathbb P(S_n=i\mid t)$ attainable by permitted controllers. They
admit the formal nonnegative representation
\begin{equation}\label{eq:v33-formal-recursion}
 \begin{split}
 q_{n,i}(t)&=\sum_{|\alpha|=n}c_{n,i,\alpha}\,(k^a(t))^\alpha,\\
 c_{n,j,\alpha}
    &=\sum_{i,k:\,\alpha_k>0}
          c_{n-1,i,\alpha-e_k}A_{n,kj}(i),\qquad
 \sum_i c_{n,i,\alpha}=\binom n\alpha.
 \end{split}
\end{equation}
Thus there are $TM^2J$ one-step entries and at most
$M\binom{n+J-1}{J-1}$ formal occupancy coefficients at time $n$.
No variable is indexed by a complete history or by a command grid.
All prior-dependent coefficients are determined by cell moments
$\mu((k^a)^\alpha)$ with $|\alpha|\le N$.
\end{theorem}
\begin{proof}
For an actual controller let $U_n(j\mid i,g,x)$ be its transition
probabilities, after averaging its fresh private randomization.
Conditional on the latent parameter and the preceding state, the
current command has law $\nu_n$. Integrating first over this command
and then over its report gives \eqref{eq:v33-q-recursion}, with
$A_n(i)$ defined by \eqref{eq:v33-realize-row}. Each such matrix
belongs to the body in Lemma~\ref{lem:v33-body}.
Conversely, that lemma supplies an admissible transition kernel for
each of the finitely many rows in \eqref{eq:v33-q-recursion}. Use
these kernels at their prescribed stages. Induction proves that their
conditional state probabilities are exactly $q_{n,i}$. These
probabilities are an offline description of a controller's law; the
controller retains only its label, not the functions $q_{n,i}$.

Multiplication by one detector cell proves the first two identities
in \eqref{eq:v33-formal-recursion}. All coefficients are nonnegative.
Summing over the new label and using $\sum_j A_{n,kj}(i)=1$ gives
$\sum_i c_{n,i,\alpha}=\sum_{k:\alpha_k>0}\binom{n-1}{\alpha-e_k}
=\binom n\alpha$. In particular $0\le q_{n,i}\le1$ and
$\sum_i q_{n,i}=1$.

Conditional on a full command-report history of positive evidence,
the latent parameter and the retained state are independent; histories
of zero evidence contribute nothing to the following expectations. To check this directly, the
joint density with respect to $\mu$, for a fixed history, is its
likelihood product multiplied by the probability of the state path
under the transition kernels. The latter factor depends on the
history and the private coins, not on the latent parameter. Summing
state paths preserves this factorization. It follows that, for every
decoder table,
\[
 \mathbb E\bigl[p^a_{nq}(H_n^{\rm hist})\,d_{nq}(S_n)\bigr]
       =\mathbb E\bigl[H^a_{nq}(t)\,d_{nq}(S_n)\bigr]
       =\sum_i b_{niq}d_{nq}(i).
\]
Here $H_n^{\rm hist}$ denotes the random history;
it does not enter the decoder. Expanding the square therefore yields
\begin{equation}\label{eq:v33-fixed-decoder}
 R_n=B_n(a)+\sum_{i,q}\omega_{nq}
          \bigl(w_{ni}d_{nq}(i)^2-2b_{niq}d_{nq}(i)\bigr).
\end{equation}
Completing each square gives the displayed centroids and
\eqref{eq:v33-exact-value}. A separate decoder coordinate is used
at each checkpoint, so these minimizations are simultaneous. The
transition matrices, in contrast, are common to every checkpoint.

The product of the finitely many implementability bodies is compact.
The quantities $w$ and $b$ are polynomial functions of its entries.
Since $0\le b_{niq}\le w_{ni}$, the ratio $b_{niq}^2/w_{ni}$ lies
between zero and $w_{ni}$ and extends continuously at $w_{ni}=0$.
Thus the objective in \eqref{eq:v33-exact-value} is continuous and
attains its minimum. The constructive converse for the matrices and
the centroid tables realize this minimum.

A future product $H^a_{nq}$ has degree at most $N-n$ in the cells. Hence $w$
and $b$ involve cell moments of degree at most $N$. For completeness,
if $L_h$ is a history likelihood of degree $n$, then
\begin{equation}\label{eq:v33-baseline}
 B_n(a)=\sum_{x_1,\ldots,x_n}\int_{G^n}\sum_q\omega_{nq}
       \frac{\mu(L_hH^a_{nq})^2}{\mu(L_h)}
               \prod_{s=1}^n\nu_s(\mathrm dg_s).
\end{equation}
Its evaluated moments have degree at most $N$. In the general
finite-cell model, set its integrand to zero whenever $\mu(L_h)=0$.
This is well defined because
\[
 0\le\mu(L_hH^a_{nq})\le\mu(L_h),\qquad
 0\le\frac{\mu(L_hH^a_{nq})^2}{\mu(L_h)}\le\mu(L_h)
\]
with that convention. These bounds justify integrability, including
impossible histories, and prove the finite-moment assertion. In the
positive monomial specialization the stronger denominator bound
$\mu(L_h)\ge\kappa^n$ is available, as used in the geometric and
approximation theorems.
The integral in \eqref{eq:v33-baseline} is a controller-independent
coefficient evaluation; the theorem does not assert a polynomial-time
procedure for evaluating it or solving the nonconvex minimum.
\end{proof}

\begin{remark}[Which tree has been removed]\label{rem:v33-tree}
The history-state variables in \eqref{eq:v32-compatible} have been
removed from the controller optimization, not hidden in an additional
retained state. The body \eqref{eq:v33-body} also integrates the
continuous command exactly. Formula~\eqref{eq:v33-baseline} can still
be costly to evaluate. The distinction matters: representation size,
coefficient evaluation and global optimization are different questions.
For fixed $J$ the number of formal occupancy coefficients grows
polynomially with the horizon, while the stated nonlinear optimization
need not be efficiently solvable. The formula concerns prescribed
acquisition and mean regret. It does not replace the worst-history
criterion by a centroid objective, nor the state-controlled criterion
by an experiment-independent offset $B_n$.
\end{remark}

\subsection{The evaluation algebra and attainable geometry}
\begin{proposition}[Formal closure and its evaluation]
\label{prop:v34-evaluation}
For fixed cells define
\[
 \mathcal F_m(k^a)=\operatorname{span}
       \{(k^a)^\alpha:|\alpha|\le m\},\qquad \mathcal F_0=\R\one.
\]
Then $\mathcal F_m\mathcal F_s\subset\mathcal F_{m+s}$, and the
same space is spanned by the products of exactly $m$ cells. Every
conditional occupancy at time $n$ belongs to $\mathcal F_n$.
If the cells span $W_A$, then
\begin{equation}\label{eq:v34-evaluation}
                 \mathcal F_m(k^a)=W_{mA}.
\end{equation}
The formal coefficients in \eqref{eq:v33-formal-recursion} are
therefore propagated before evaluation in a function space whose
linear relations may depend on the experiment.
\end{proposition}
\begin{proof}
Multiplication adds degrees. Multiplying a product of degree $d<m$
by $(\sum_k k_k^a)^{m-d}=1$ writes it as a sum of degree-$m$
products. The occupancy assertion follows by induction in
\eqref{eq:v33-q-recursion}. If the cells span $W_A$, each cell
product is a linear combination of monomials in $mA$. Conversely,
write each factor $t^{a_i}$ as a linear combination of cells and
multiply $m$ such identities. Padding by constants covers shorter
products. This proves both inclusions in \eqref{eq:v34-evaluation}.
\end{proof}

The equality identifies the future test space, not the dimension or
the measure of its attained prediction image. The differential of the
normalized product must still be paired with this space. In the
monomial experiment, Lemma~\ref{lem:binomial-tangent} and the positive
pairing give the acquired dimension, and
Lemma~\ref{lem:newton-attainment} supplies its collision-uniform
lower mass. None of those conclusions is inferred from the number
of formal occupancy coefficients.

\begin{proposition}[A two-cell realization]
\label{prop:v34-two-cell}
Let $B$ be a Borel subset of $\Theta$ with $p=\mu(B)\in(0,1)$ and
put $k_1=\tfrac14+\tfrac12\one_B$, $k_2=1-k_1$. Use independent
uniform acquisition commands on $G$. For every $m\ge1$,
\[
            \mathcal F_m=\operatorname{span}\{\one,\one_B\}.
\]
Every acquired future prediction is determined by the single posterior
weight of $B$. For every $n\ge1$, the all-failure prediction image
contains a nondegenerate interval of these weights. A future failure
query with unequal command coordinates distinguishes them. Thus this
positive finite-cell experiment has one-dimensional attained
prediction geometry at every nontrivial checkpoint whose menu contains
such a distinguishing query, although its
formal degree-$n$ two-cell representation has $n+1$ coefficients
per label. Taking $\Theta=[0,1]$ and $B=[0,1/2]$ gives a
nonmonomial example.
\end{proposition}
\begin{proof}
Every cell and cell product is constant on each of $B$ and $B^c$,
while $\one_B=2k_1-\tfrac12$, proving the space assertion. Write
$z$ for the posterior weight of $B$. Every future test in this space
has posterior expectation $zH|_B+(1-z)H|_{B^c}$.

Take all but one of the $n$ commands on the interior diagonal
$g_1=g_2=c$, $\eta<c<1-\eta$, and consider the all-failure word.
Those other factors are constant. At the remaining command the
failure likelihoods on $B$ and $B^c$ are, respectively,
\[
 L_B=1-\frac{3g_1+g_2}{4},\qquad
 L_{B^c}=1-\frac{g_1+3g_2}{4}.
\]
Consequently
\[
 z=\frac{pL_B}{pL_B+(1-p)L_{B^c}},\qquad
 \left.\frac{\partial z}{\partial g_1}\right|_{g_1=g_2=c}
       =-\frac{p(1-p)}{2(1-c)}\ne0.
\]
An interval of weights is attained by interior commands. A future
failure factor with $g_1\ne g_2$ has different values on $B$ and
$B^c$, so distinguishes them. Additional required future trials can
use constant diagonal failure factors. This also proves the claim
when the horizon counts every one of those trials.
\end{proof}

\subsection{Exact removal of private transition randomization}
We record the finite-action purification argument to specify its use
here. Its abstract mechanism is classical \cite{DWW1951}; the relevant
atomless measure in this application is the command law, not the prior.

\begin{lemma}[Finite-moment purification]\label{lem:v33-purification}
Let $\nu$ be an atomless Borel probability on $G$, let
$\phi_1,\ldots,\phi_d\in L^1(\nu)$, and let
$u:G\to\Delta_M$ be Borel. There is a Borel partition
$E_1,\ldots,E_M$ such that
\[
 \int_{E_j}\phi_\ell\,\mathrm d\nu
       =\int_G\phi_\ell u_j\,\mathrm d\nu
              \qquad(1\le j\le M,\ 1\le\ell\le d).
\]
\end{lemma}
\begin{proof}
In the weak-star compact set of simplex-valued elements of
$L^\infty(\nu)^M$, impose the finitely many displayed equalities.
The resulting nonempty compact convex set has an extreme point.
Suppose two coordinates of that point are both at least $\epsilon>0$
on a measurable set $E$ of positive measure. Such an $E$, pair of
coordinates and $\epsilon$ exist if the point fails to be a simplex
vertex on a set of positive measure. Atomlessness splits $E$ into
$d+1$ disjoint sets $E_s$ of positive measure. The $d+1$ vectors
$(\int_{E_s}\phi_\ell\,\mathrm d\nu)_{\ell=1}^d$ are linearly
dependent. A nonzero linear combination of their indicators gives
a bounded function $h$, supported in $E$, with
$\int h\phi_\ell\,\mathrm d\nu=0$ for every $\ell$. Rescale so
$|h|\le\epsilon$. Add $h$ to the first of the two coordinates and
subtract it from the second, or make the opposite changes. Both
resulting functions are simplex-valued and satisfy every constraint;
they are distinct and have the proposed extreme point as their mean.
This contradiction proves that the extreme point is a vertex almost
everywhere. Its coordinate sets, modified on a null set, give a Borel
partition with the required integrals.
\end{proof}

\begin{theorem}[Deterministic realization of the entire mean-risk vector]
\label{thm:v33-pure}
Suppose every prescribed command law $\nu_n$ is atomless. Every
$M$-label controller using fresh private transition coins has a
$M$-label deterministic Borel replacement with the same decoder
and the same conditional state probabilities $q_{n,i}(t)$ at every
time and every latent parameter. For fixed transition kernels, the replacement is independent of the
calibration and preserves these laws at every calibration with the
cell index set and observation channel fixed. For calibration-indexed kernels the assertion
is pointwise in calibration. In particular,
it has exactly the same vector of unconditional checkpoint regrets.
Consequently the minimum in \eqref{eq:v33-exact-value} is attained
without transition randomization.
\end{theorem}
\begin{proof}
For a fixed time and preceding state, apply
Lemma~\ref{lem:v33-purification} to the accepted-report kernel for
report $k$ with weight $g_k$, separately for each $k$. Apply it to
the failure kernel with the $J$ weights $1-g_k$. The resulting
deterministic rules preserve every entry in
\eqref{eq:v33-realize-row}. Repeat for the finitely many times and
preceding states. Recursion~\eqref{eq:v33-q-recursion} then preserves
all the conditional state probabilities. The weights used in
purification contain neither $a$ nor $t$, so the same replacement
preserves those probabilities at every calibration when both programs
use the same fixed kernels. This asserts no jointly measurable
purification of an arbitrary calibration-indexed family. Formula
\eqref{eq:v33-fixed-decoder} proves equality of the mean-risk
vectors for the original decoder tables. Apply this replacement to
a minimizer in Theorem~\ref{thm:v33-moment-program}.
\end{proof}

This theorem does not purify the risk at each individual history. The
replacement can assign different labels to a particular command-report
prefix. It therefore implies no equality of the two worst-history
optima. Nor does it remove a persistent public seed that also indexes
a decoder program; that is a different resource. A singular full-support
latent prior is allowed because the atomlessness used in the proof
belongs entirely to acquisition.

\subsection{Polyhedral transition rules for scalarized regret}
A more explicit description is available for a weighted sum of mean
regrets. It is important not to infer a convex minimax identity for
\eqref{eq:v33-exact-value} from this description.

\begin{theorem}[Polyhedral rows]\label{thm:v33-polyhedral}
Fix nonnegative weights $\lambda_1,\ldots,\lambda_T$. In the
prescribed-command model there is a globally optimal deterministic
controller for $\sum_n\lambda_n R_n$ with the following form at each
time and preceding state. After an accepted report $k$ its new label
is independent of the command. After failure its new label minimizes
one of $M$ affine functions of $g$. Thus its failure regions are intersections of the command cube with
at most $M$ Borel polyhedral regions, using the half-open tie
convention specified in the proof.
The assertion holds for arbitrary prescribed command laws, including
atomic ones.
\end{theorem}
\begin{proof}
Compactness of the implementability bodies and decoder cube gives a
global minimizer of the fixed-decoder expression
\eqref{eq:v33-fixed-decoder}, summed with weights $\lambda_n$.
Freeze its decoder tables. Freeze all transition matrices except
$A_s(i)$, for one time $s$ and preceding state $i$. The resulting
objective is affine in this matrix. Indeed, a path uses time $s$
only once, so induction in \eqref{eq:v33-q-recursion} makes every
later occupancy linear in these entries. Write its nonconstant part
as $\sum_{k,j}C_{kj}A_{s,kj}(i)$.

The minimum support formula in Lemma~\ref{lem:v33-body}, obtained
by replacing $C$ by $-C$, is attained as follows: after report $k$
choose a fixed index minimizing $C_{kj}$; after failure choose the
smallest index minimizing $\sum_k(1-g_k)C_{kj}$. These are permitted
deterministic rules. Writing these affine functions as $F_j$, the precise region for
label $j$ is
\begin{equation}\label{eq:v34-ties}
 P_j=G\cap\bigcap_{l<j}\{F_j<F_l\}
         \cap\bigcap_{l>j}\{F_j\le F_l\}.
\end{equation}
These regions form a Borel partition, possibly with half-open faces.
The tie convention is retained even when $\nu$ gives a boundary
positive mass. Replacing the one row by this
minimizer cannot increase the objective. Since the original value
was already globally minimal, its value remains unchanged.

Perform this replacement successively for all time-state pairs.
At each step the full controller and the frozen decoders remain
globally optimal, and previously replaced rules retain the asserted
form. After finitely many replacements the required controller has
been obtained. No atomlessness, minimax exchange or convexity of the
joint risk set was used.
\end{proof}

The theorem identifies the geometry of an optimal transition rule,
not its numerical coefficients. The coefficient matrices depend on
the other transitions and the decoder. In particular, the theorem
does not assert that a weighted optimum solves the maximum-checkpoint
problem for some weights. That problem is given exactly by
\eqref{eq:v33-exact-value} and, for atomless commands, by
Theorem~\ref{thm:v33-pure}.

\subsection{Collision limits in the same controller space}
We now return to the positive monomial experiment on $[0,1]$ and
the compact exponent chamber of Theorem~\ref{thm:resolution-main}.
The general finite-cell theorem identifies its controller domain;
the following continuity assertion additionally uses this compact,
uniformly positive calibration family.
\begin{theorem}[Compact collision transfer]\label{thm:v33-collision}
Fix the detector matrix, prior, horizon and prescribed command laws.
Let the known calibration vary in the compact chamber $K$, and use
the fixed attainable failure-factor query basis of
Theorem~\ref{thm:resolution-main}. The implementability domain in
\eqref{eq:v33-exact-value} is independent of $a$, and its objective
is jointly continuous in $a$ and the one-step matrices. Its value
is continuous on $K$.

If $a_s\to a_0$ and $A^{(s)}$ are optimal matrices, a subsequence
converges to optimal matrices at $a_0$. Along that subsequence their
conditional state probabilities converge uniformly in $t\in[0,1]$
at every stage. These assertions include all additive collisions.
For uniform commands the exact value in this same domain satisfies
\[
 c\Xi_N(M,a)\le
    \min_A\max_n\left(B_n-\sum_{i,q}\omega_{nq}b_{niq}^2/w_{ni}\right)
           \le C\Xi_N(M,a)
\]
with the constants and all-budget collision profile of
Theorem~\ref{thm:resolution-main}.
\end{theorem}
\begin{proof}
The monomials defining the cells are jointly continuous in $(a,t)$
on $K\times[0,1]$; positive one-step exponents are bounded away
from zero there. Consequently all cell products and fixed future
queries are uniformly continuous on this compact set. Their moments
are continuous for the fixed prior, without a density hypothesis.
The finite polynomial recursion proves joint continuity of $q,w,b$.
The positive denominator in \eqref{eq:v33-baseline} proves continuity
of $B_n$ by bounded convergence, uniformly on the finite-dimensional
compact domain. At zero occupancy the bound $b^2/w\le w$ again
proves continuity of the perspective term. Thus the entire objective
is continuous on $K$ times one compact, calibration-independent set.

Let $A_0$ minimize the objective at $a_0$. Substitution of $A_0$
at $a_s$ gives the upper limit inequality for the values. For the
lower limit inequality, take convergent subsequences of minimizers
in the fixed compact domain and use joint continuity. The same
argument proves optimality of every such limit. Uniform convergence
of the state probabilities follows from the finite polynomial
recursion and uniform convergence of the cells.

The last assertion is Theorem~\ref{thm:resolution-main} applied to
the exact equality in Theorem~\ref{thm:v33-moment-program}. It uses
the complete acquired collision profile, rather than only its
one-dimensional floor.
\end{proof}

At an additive collision, equal exponent sums in the evaluation of
$(k^a)^\alpha$ are simply collected. Formula
\eqref{eq:v33-formal-recursion} is used before this quotient and
never divides by an exponent gap. Thus the changing dimension of the
posterior test space does not require a changing space of controller
coefficients. The continuity conclusion is at fixed $M$; the displayed
classification is the separate statement uniform in $M$ and $a$.
Neither conclusion specifies a sharp leading constant for arbitrary
$M$, and neither is inferred from the precision chosen for a
certificate.
